\documentclass{article}
\usepackage{amsmath, amssymb}
\usepackage[margin=1in]{geometry}

\title{HW8}
\author{
Greyson Knippel \\[4pt]
MATH 339 --- Gonzaga University \\[4pt]
compiled with \LaTeX{}
}
\date{\today}

\begin{document}
\maketitle

\section*{4.1 - 17}

\[
T\left(\begin{bmatrix} x \\ y \end{bmatrix}\right) = \begin{bmatrix} -x \\ y \end{bmatrix},
\qquad
\mathbf{u} = \begin{bmatrix} -2 \\ 3 \end{bmatrix},
\qquad
\mathbf{v} = \begin{bmatrix} 2 \\ -2 \end{bmatrix}.
\]

\subsection*{a.}

T negates the first component and keeps the second, so:
\[
T(\mathbf{u}) = \begin{bmatrix} 2 \\ 3 \end{bmatrix},
\qquad
T(\mathbf{v}) = \begin{bmatrix} -2 \\ -2 \end{bmatrix}.
\]

\subsection*{b.}

\[
\mathbf{u} + \mathbf{v} = \begin{bmatrix} 0 \\ 1 \end{bmatrix}
\quad \Rightarrow \quad
T(\mathbf{u} + \mathbf{v}) = \begin{bmatrix} 0 \\ 1 \end{bmatrix},
\]
\[
T(\mathbf{u}) + T(\mathbf{v})
= \begin{bmatrix} 2 \\ 3 \end{bmatrix} + \begin{bmatrix} -2 \\ -2 \end{bmatrix}
= \begin{bmatrix} 0 \\ 1 \end{bmatrix}
= T(\mathbf{u} + \mathbf{v}).
\]
so yes, $T(\mathbf{u} + \mathbf{v}) = T(\mathbf{u}) + T(\mathbf{v})$

\subsection*{c.}

one example working isn't enough, so we check both properties for vectors
$\begin{bmatrix} x_1 \\ y_1 \end{bmatrix}, \begin{bmatrix} x_2 \\ y_2 \end{bmatrix}$
and scalar $c$:
\[
T\left(\begin{bmatrix} x_1 + x_2 \\ y_1 + y_2 \end{bmatrix}\right)
= \begin{bmatrix} -x_1 - x_2 \\ y_1 + y_2 \end{bmatrix}
= \begin{bmatrix} -x_1 \\ y_1 \end{bmatrix} + \begin{bmatrix} -x_2 \\ y_2 \end{bmatrix}
= T\left(\begin{bmatrix} x_1 \\ y_1 \end{bmatrix}\right) + T\left(\begin{bmatrix} x_2 \\ y_2 \end{bmatrix}\right),
\]
\[
T\left(c\begin{bmatrix} x \\ y \end{bmatrix}\right)
= \begin{bmatrix} -cx \\ cy \end{bmatrix}
= c\begin{bmatrix} -x \\ y \end{bmatrix}
= c\,T\left(\begin{bmatrix} x \\ y \end{bmatrix}\right),
\]
so yes, T is a linear transformation because its a reflection across the y axis.

\section*{4.1 - 27}

\[
T(x^2) = 2x - 1, \qquad T(-3x) = x^2 - 1.
\]


since T is linear, we can pull the scalar out of $T(-3x)$:
\[
T(-3x) = -3\,T(x) = x^2 - 1
\quad \Rightarrow \quad
T(x) = -\tfrac{1}{3}(x^2 - 1) = -\tfrac{1}{3}x^2 + \tfrac{1}{3}.
\]
but we still know nothing about $T(1)$

\subsection*{a.}

using linearity,
\[
T(2x^2 - 3x + 2) = 2\,T(x^2) - 3\,T(x) + 2\,T(1).
\]
this needs $T(1)$, and $1$ is not a linear combination of $x^2$ and $x$, so the given
information doesn't determine it. so no, $T(2x^2 - 3x + 2)$ cannot be determined

\subsection*{b.}

$3x^2 - 4x$ only uses $x^2$ and $x$, so we can write it in terms of the given inputs:
\[
3x^2 - 4x = 3(x^2) + \tfrac{4}{3}(-3x).
\]
then by linearity,
\[
T(3x^2 - 4x) = 3\,T(x^2) + \tfrac{4}{3}\,T(-3x)
= 3(2x - 1) + \tfrac{4}{3}(x^2 - 1)
= 6x - 3 + \tfrac{4}{3}x^2 - \tfrac{4}{3}
= \tfrac{4}{3}x^2 + 6x - \tfrac{13}{3}.
\]
so yes, $T(3x^2 - 4x) = \tfrac{4}{3}x^2 + 6x - \tfrac{13}{3}$.

\end{document}
